\begin{lemma}
Let $U\subset\mathbb{R}^n$ be open, let $\sigma=(\sigma^1,\dots,\sigma^m)$ be smooth vector fields
$\sigma^i:\mathbb{R}^n\to\mathbb{R}^n$ and let $a:\mathbb{R}^m\to\mathbb{R}^{m\times m}$ be such
that $a(\xi)$ is positive semi-definite (hence symmetric) for every
$\xi\in\mathbb{R}^m\setminus\{0\}$. Let $u,v:\mathbb{R}^n\to\mathbb{R}$ be of class $C^2$ on $U$
and let $x_0\in U$ be a local maximum point of $u-v$ with $\mathfrak{X}u(x_0)\neq0$. Then
$$\mathfrak{X}v(x_0)=\mathfrak{X}u(x_0),\qquad
\mathfrak{X}\mathfrak{X}^\star v(x_0)-\mathfrak{X}\mathfrak{X}^\star u(x_0)
\ \text{ is positive semi-definite},\qquad
\mathcal{L}v(x_0)\leq\mathcal{L}u(x_0),$$
the middle assertion being $\mathfrak{X}\mathfrak{X}^\star
u(x_0)\leq\mathfrak{X}\mathfrak{X}^\star v(x_0)$ in the order of Section~2.1 of the paper. This is
the first assertion of Lemma \texttt{lem:vweak} of the paper.
\end{lemma}
\begin{proof}
Since $U$ is open and $x_0\in U$, and $u,v$ are of class $C^2$ on $U$, both $u$ and $v$ are of
class $C^2$ in a neighbourhood of $x_0$ (namely on $U$ itself); in particular both are
differentiable at $x_0$. Each $\sigma^i$ is smooth, hence of class $C^1$ in a neighbourhood of
$x_0$.

\emph{Step 1: $\mathfrak{X}v(x_0)=\mathfrak{X}u(x_0)$.} Apply \noderef{HGradEqAtLocalMax} to
$\sigma$, $u$, $v$ and $x_0$: its hypotheses are that $u$ and $v$ are differentiable at $x_0$ and
that $x_0$ is a local maximum point of $u-v$, all of which hold. It gives
$\mathfrak{X}u(x_0)=\mathfrak{X}v(x_0)$. Write $\xi_0:=\mathfrak{X}u(x_0)=\mathfrak{X}v(x_0)$; by
hypothesis $\xi_0\neq0$.

\emph{Step 2: $Du(x_0)=Dv(x_0)$.} Apply \noderef{FDerivEqAtLocalMaxSub} to $u$, $v$ and $x_0$: its
hypotheses are that $u$ and $v$ are differentiable at $x_0$ and that $x_0$ is a local maximum point
of $u-v$, all of which hold. It gives $Du(x_0)=Dv(x_0)$.

\emph{Step 3: $D^2u(x_0)[w,w]\leq D^2v(x_0)[w,w]$ for every $w\in\mathbb{R}^n$.} Apply
\noderef{SndFDerivLeAtLocalMaxSub} to $U$, $u$, $v$ and $x_0$: its hypotheses are that $U$ is open,
that $x_0\in U$, that $u$ and $v$ are of class $C^2$ on $U$ and that $x_0$ is a local maximum point
of $u-v$, all of which hold by hypothesis. It gives exactly the claim.

\emph{Step 4: the symmetrized horizontal second derivatives.} Steps 2 and 3 are exactly the two
hypotheses of \noderef{HGradHHessComparison}, whose remaining hypotheses ($u,v$ of class $C^2$
near $x_0$, each $\sigma^i$ of class $C^1$ near $x_0$) also hold. It yields that
$$P:=\mathfrak{X}\mathfrak{X}^\star v(x_0)-\mathfrak{X}\mathfrak{X}^\star u(x_0)$$
is positive semi-definite, which is the second assertion.

\emph{Step 5: the trace comparison.} Put $A:=a(\xi_0)$. Since $\xi_0\neq0$, $A$ is positive
semi-definite by hypothesis; in particular $A$ is symmetric. Apply \noderef{TraceMonotone} with
this $A$, with $X:=\mathfrak{X}\mathfrak{X}^\star v(x_0)$ and
$Y:=\mathfrak{X}\mathfrak{X}^\star u(x_0)$: its hypotheses are that $A$ is positive semi-definite
(just checked) and that $X-Y=P$ is positive semi-definite (Step 4). It gives
\begin{equation*}\operatorname{Tr}\Big(A\,\mathfrak{X}\mathfrak{X}^\star u(x_0)\Big)\ \leq\
\operatorname{Tr}\Big(A\,\mathfrak{X}\mathfrak{X}^\star v(x_0)\Big). \tag{2}\end{equation*}

\emph{Step 6: symmetrization does not change the trace against a symmetric matrix.} The matrix
$A=a(\xi_0)$ is symmetric (Step 5), so \noderef{TraceMulSymmetrizeEq} applies to $A$ and to every
$B\in\mathbb{R}^{m\times m}$ and gives
$\operatorname{Tr}\big(A\cdot\tfrac12(B+B^T)\big)=\operatorname{Tr}(AB)$. Taking
$B=\mathfrak{X}\mathfrak{X}u(x_0)$ and then $B=\mathfrak{X}\mathfrak{X}v(x_0)$, and recalling that
by the definition of the symmetrization (\noderef{HHessSym}) one has
$\mathfrak{X}\mathfrak{X}^\star w(x_0)=\tfrac12\big(\mathfrak{X}\mathfrak{X}w(x_0)
+\mathfrak{X}\mathfrak{X}w(x_0)^T\big)$,
\begin{equation*}\operatorname{Tr}\Big(A\,\mathfrak{X}\mathfrak{X}^\star u(x_0)\Big)
=\operatorname{Tr}\Big(A\,\mathfrak{X}\mathfrak{X}u(x_0)\Big),\qquad
\operatorname{Tr}\Big(A\,\mathfrak{X}\mathfrak{X}^\star v(x_0)\Big)
=\operatorname{Tr}\Big(A\,\mathfrak{X}\mathfrak{X}v(x_0)\Big). \tag{3}\end{equation*}
This is the paper's remark, just after \texttt{eq:maineq}, that the equation can equivalently be
written with the symmetrized horizontal second derivative.

\emph{Step 7: conclusion.} By the definition of the operator (\noderef{SubellipticOp}) and
$\mathfrak{X}u(x_0)=\mathfrak{X}v(x_0)=\xi_0$,
$$\mathcal{L}u(x_0)=-\operatorname{Tr}\Big(A\,\mathfrak{X}\mathfrak{X}u(x_0)\Big),\qquad
\mathcal{L}v(x_0)=-\operatorname{Tr}\Big(A\,\mathfrak{X}\mathfrak{X}v(x_0)\Big).$$
Combining this with (3) and then (2),
$$-\mathcal{L}u(x_0)=\operatorname{Tr}\Big(A\,\mathfrak{X}\mathfrak{X}^\star u(x_0)\Big)
\ \leq\ \operatorname{Tr}\Big(A\,\mathfrak{X}\mathfrak{X}^\star v(x_0)\Big)=-\mathcal{L}v(x_0),$$
that is, $\mathcal{L}v(x_0)\leq\mathcal{L}u(x_0)$.
\end{proof}
